% Source: https://artofproblemsolving.com/community/c3h2755249_denesting_radicals_harder
% Stable identity: aops-topic-2755249-q01
\textbf{(a)} Try $\left(\sqrt p+\sqrt q+\sqrt r\right)^2=p+q+r+2\sqrt{pq}+2\sqrt{pr}+2\sqrt{qr}$ with $pq=6$, $pr=10$, $qr=15$. Then $(pqr)^2=900$, so $pqr=30$ and $p=2$, $q=3$, $r=5$; indeed $p+q+r=10$. Hence
\[\sqrt{10+2\sqrt6+2\sqrt{10}+2\sqrt{15}}=\boxed{\sqrt2+\sqrt3+\sqrt5}.\]

\textbf{(b)} The term $\sqrt{30}=\sqrt2\sqrt3\sqrt5$ suggests the same $s=\sqrt2+\sqrt3+\sqrt5$. Using (a),
\[s^3=s\cdot s^2=\left(\sqrt2+\sqrt3+\sqrt5\right)\left(10+2\sqrt6+2\sqrt{10}+2\sqrt{15}\right).\]
Multiplying out and collecting terms:
\begin{itemize}
\item $\sqrt2$: $10+6+10=26$ (from $\sqrt2\cdot10$, $\sqrt3\cdot2\sqrt6$, $\sqrt5\cdot2\sqrt{10}$);
\item $\sqrt3$: $4+10+10=24$ (from $\sqrt2\cdot2\sqrt6$, $\sqrt3\cdot10$, $\sqrt5\cdot2\sqrt{15}$);
\item $\sqrt5$: $4+6+10=20$ (from $\sqrt2\cdot2\sqrt{10}$, $\sqrt3\cdot2\sqrt{15}$, $\sqrt5\cdot10$);
\item $\sqrt{30}$: $2+2+2=6$.
\end{itemize}
So $s^3=26\sqrt2+24\sqrt3+20\sqrt5+6\sqrt{30}$, and the cube root is $\boxed{\sqrt2+\sqrt3+\sqrt5}$.

\textbf{(c)} Since $108=2^2\cdot3^3$, we have $\sqrt[6]{108}=\sqrt[6]{2^2}\cdot\sqrt[6]{3^3}=\sqrt[3]2\,\sqrt3$, and $\sqrt[3]4=\left(\sqrt[3]2\right)^2$. So
\[3+\sqrt[3]4+2\sqrt[6]{108}=\left(\sqrt3\right)^2+\left(\sqrt[3]2\right)^2+2\sqrt3\,\sqrt[3]2=\left(\sqrt3+\sqrt[3]2\right)^2,\]
and the expression equals $\boxed{\sqrt3+\sqrt[3]2}$.
